\fancyhead[L]{\small MAGNETIC OBSERVATIONS.}
\SourceLeaf{523}\pdfbookmark[0]{Northern Lights on 7 January 1831}{nordlicht_1831_januar}
\begin{center}
\hrule\vspace{2mm}
Göttingische gelehrte Anzeigen. 1831. February 28.
\vspace{2mm}\hrule\vspace{5mm}
\end{center}
Professor \textsc{Gerling} in Marburg has submitted a note to the Royal Society on his observation
\begin{center}\itshape of the northern lights seen on 7 January this year\end{center}
\noindent which broadly agrees with reports already received from other places, but also touches on a point that deserves particular attention and seems to have received too little notice so far. We therefore give an extract from it here.

The phenomenon had been seen in Marburg from as early as 6 o'clock. Mr \textsc{Gerling}, however, was informed of it only at 8 o'clock, and by then there was nothing unusual to be seen anywhere in the northern sky, down to the lowest point visible from the windows of his home. Around 9 o'clock, however, conspicuous red streaks appeared again in the northern sky, and Mr \textsc{Gerling} immediately went to the Schlossberg, which commands an unobstructed view, so that he could still see as much of the phenomenon as possible.

At first only red streaks and patches were visible in the sky, across a span of about 50—60 degrees between northeast and northwest. They extended discontinuously along the stated arc in azimuth and reached up to about 45 degrees in elevation on average. In the middle of that arc in azi\SourceLeaf{524}muth, around the meridian and across an estimated span of about 30—40 degrees in azimuth, black patches appeared in the otherwise clear sky. In appearance they resembled nothing but small black clouds. These patches gradually increased in number and finally merged to form the dark segment that, by all accounts, seems characteristic of the northern lights. At the same time the red patches mentioned first grew more intense and arranged themselves in rays pointing towards the black segment. White and yellowish rays then also appeared from it between the red rays. They may have reached an elevation of about 50 degrees without any strikingly sudden upward shooting.

‘So far, Mr \textsc{Gerling} continues, this observation seems to agree entirely with what others saw at the same time and on earlier occasions of northern lights. It would therefore scarcely be worth mentioning if I had not noticed something that, to my knowledge, has come up neither in connection with this display nor anywhere else, as far as I have been able to find. The stars of Cygnus across which the intense white and red rays passed, and the star $\alpha$ in Lyra, \textit{which lay deep within the black segment}, all remained just as visible and appeared just as bright, with no perceptible loss at all. On the puzzling question of what exactly the dark segment is, this fact seems to yield at least the negative conclusion that it is \textit{no ordinary cloud}, since such a cloud could not let starlight through.’

As early as the northern lights of 22 October 1804, \textsc{Wrede} remarked, without giving this reason, that the dark segment was \textit{wrongly} called a cloud. \textsc{Gilbert} defends the expression and adds that he noticed nothing in the dark segment that would have led him to think that what he saw there was anything other than a dark cloud. The opinion of \textsc{Mayer} in his handbook of physical astronomy, that the denser air at the horizon, laden with vapours, suffices to explain the dark segment, also seems hard to reconcile with the fact noted by Mr \textsc{Gerling}.

Mr \textsc{Gerling} adds that in the earlier hours, when the segment, which varied in extent, reached very high, a credi\SourceLeaf{525}ble witness had seen the star $\alpha$ in Lyra shining \textit{within the segment} as brightly as at any other time. Another, at a time when the dark segment had yet to reach that star, had seen \textit{other stars within the segment}.

Mr \textsc{Gerling} also attached an extract from his meteorological journal for 5—9 January. It shows nothing remarkable apart from a rise in the barometer of three quarters of an inch from the afternoon of 6 January to the evening of 7 January. On 7 January the wind was from the north.

The observations that Professor \textsc{Harding} made here in Göttingen of these northern lights agree in essentials with those reported from other places. It is worth mentioning, however, that while the phenomenon lasted the magnetic needle moved about three quarters of a degree northwards from its usual position and had returned to that position by the following morning.

\bigskip
